\documentclass[11pt,a4paper]{article}
\usepackage[margin=0.85in,top=1in,bottom=1in]{geometry}
\usepackage{amsmath,amssymb,amsfonts}
\usepackage{graphicx}
\usepackage{float}
\usepackage{enumitem}
\usepackage{microtype}
\usepackage{caption}
\usepackage[hidelinks]{hyperref}
\usepackage[most,breakable]{tcolorbox}
\tcbuselibrary{breakable,skins}
\usepackage{fontspec}
\setmainfont{DejaVu Serif}
\setsansfont{DejaVu Sans}
\setmonofont{DejaVu Sans Mono}
\captionsetup{font=small,labelfont=bf,skip=4pt}
\setlist{leftmargin=1.5em,topsep=2pt}

%% Breakable problem card — prevents content cutoff across pages
\newtcolorbox{solcard}[3]{
  breakable, enhanced,
  colback=white, colframe=gray!50, arc=2pt,
  title={\textbf{#1}\hfill\textsf{\small #3\,pts}},
  fonttitle=\normalsize\sffamily\bfseries,
  before upper={\small\textit{#2}\par\smallskip\hrule height 0.4pt\smallskip},
  top=4pt, bottom=6pt, left=7pt, right=7pt,
  parbox=false,
}

\title{\textbf{Flight to the Moon --- Solution}}
\author{\normalsize Y26 (2026) — English translation}
\date{}
\begin{document}\maketitle\vspace{-0.5em}
\begin{solcard}{A1}{Write the expressions for $\ddot y$ through $\dot{y}$, $k$, $m_0$, $\mu$, $u$, $g$ and $t$.}{0.50}

Let us consider the change in the momentum of the system. Let at some moment the mass of the rocket be $m$ and its speed $\dot y$. After a short time $dt$ the mass of the rocket became less by $\mu dt$, and the speed of the rocket increased by $\ddot y dt$. In this case, the combustion products in the Earth’s reference frame have a speed $u-\dot y$ directed downward. In this case, two external forces act on the system: air resistance and the gravity of the Earth, therefore the law of change in momentum $$m\dot y -mgdt-k\dot y\left|\dot y\right|dt= (m-\mu dt)(\dot y + \ddot y dt) - \mu dt(u-\dot y).$$Let us express $\ddot y$ $$-mg-k\dot y \left|\dot y\right| = m\ddot y - \mu u.$$$$\ddot y = \frac{\mu u}{m}-g-\frac{k\dot y \left|\dot y\right|}{m}$$It remains to express it through the required quantities.

\par\smallskip\noindent\textbf{Answer:} \textbackslash{}textbf{Answer:} $$\ddot y = \frac{\mu u}{m_0-\mu t}-g-\frac{k\dot y \left|\dot y\right|}{m_0-\mu t}$$
\end{solcard}

\begin{solcard}{A2}{Neglecting the drag force, find the speed of the rocket $v(m)$ and the coordinate $y(m)$ at the moment when its mass is $m$. Express your answer in terms of $m$, $g$, $u$, $\mu$ and $m_0$.}{1.00}

Let's use the answer from the previous paragraph and integrate twice, neglecting the resistance force, substituting $t=\dfrac{m_0-m}{\mu}$ $$\ddot y dt=\dfrac{\mu u}{m} dt-g dt,$$$$\int\limits_0^v d\dot y =- \int\limits_{m_0}^m u\dfrac{dm}{m} - \int\limits_0^t g dt,$$


Integrating the expression for speed over time $$\int \limits_0^y dy=-u\int \limits_0^t \ln{\left(1-\dfrac{\mu t}{m_0}\right) dt} -\int \limits_0^t gt dt.$$For интеграла логарифма известно $\int \ln{x}dx=x\ln{x}-x+C$ $$y=\dfrac{um_0}{\mu}\left(\left(1-\dfrac{\mu t}{m_0}\right)\ln{\left(1-\dfrac{\mu t}{m_0}\right)} - \left(1-\dfrac{\mu t}{m_0}\right)+1\right)-\dfrac{gt^2}{2},$$

\par\smallskip\noindent\textbf{Answer:} \textbackslash{}textbf{Answer:} $$v=-u\ln{\dfrac{m}{m_0}}-g\dfrac{m_0-m}{\mu}$$
\end{solcard}

\begin{solcard}{A3}{Using the result of the previous paragraph and assuming that $k$ weakly depends on height, find numerically at what height $H_c$ the rocket will be at the moment when the drag force becomes equal to the force of gravity.}{0.80}

Let us write down the condition for the equality of gravity and the resistance force $$kv^2=mg,$$$$k\left(-u\ln{\dfrac{m}{m_0}}-g\dfrac{m_0-m}{\mu}\right)^2=mg.$$Решая численно equation относительно $m$, we get $m=209~\text{т}$. Substituting массу в формулу for высоты подъема $y$.

\par\smallskip\noindent\textbf{Answer:} \textbackslash{}textbf{Answer:} $$H_c=124~\text{km}$$
\end{solcard}

\begin{solcard}{A4}{Considering a vertical take-off, numerically estimate the ratio $\beta_1$ of the initial to the final mass of the rocket, at which the rocket develops a speed $v_a$ (the speed of movement in a circular orbit of radius $a$).}{0.70}

For motion in an orbit with radius $a$ we have $v_a=\sqrt{\dfrac{GM_e}{a}}$, we substitute the expression for the speed $v$ from point A2, we get a numerical answer.

\par\smallskip\noindent\textbf{Answer:} \textbackslash{}textbf{Answer:} $$\beta_1=8.03$$
\end{solcard}

\begin{solcard}{B1}{Consider a point on the Earth-Moon straight line between the Earth and the Moon, at which the gravitational forces of the Earth and the Moon are equal in magnitude. Find the distance $r_c$ from the Moon to this point. Express your answer in terms of $r_m$, $M_e$ and $M_m$.}{0.30}

Equality of power at a distance $r_c$ $$\dfrac{GM_m}{r_c^2} = \dfrac{GM_e}{(r_m - r_c)^2}$$ Преобсуя we get $$r_c = \dfrac{r_m}{\sqrt{\frac{M_e}{M_m}} + 1}$$


$$\dfrac{GM_m}{r_c^2} = \dfrac{GM_e}{(r_m - r_c)^2}$$


Transforming we get


$$r_c = \dfrac{r_m}{\sqrt{\frac{M_e}{M_m}} + 1}$$

\par\smallskip\noindent\textbf{Answer:} \textbackslash{}textbf{Answer:} $$r_c = \dfrac{r_m}{1 + \sqrt{\frac{M_e}{M_m}}}=3.86\cdot10^7~\text{m}$$
\end{solcard}

\begin{solcard}{B2}{Considering the ship in the CO of the Moon, write down the expression for the specific angular momentum $L_0^m$ relative to the center of the Moon and the speed $v_2$ at the moment when the distance to the Moon is $r_c$. Express your answer using $s$, $s'$, $c$, $v_1$, $v_m$ and $r_c$.}{0.50}

The vector $\vec{v_2}$ is obtained from the velocity vector $\vec{v_1}$ by subtracting the vector $\vec{v_m}$, from which a velocity triangle can be formed. Using the cosine theorem $$v_2 = \sqrt{v_1^2+v_m^2-2v_mv_1\sin\alpha}.$$


Let us write the specific angular momentum by definition $$\vec{L_0^m}=\left[\vec{r_c},\vec{v_1}-\vec{v_m}\right]=\left[\vec{r_c},\vec{v_1}\right]- \left[\vec{r_c},\vec{v_m}\right].$$Возьмем magnitude выражения $$L_0^m = r_c\left(v_1\sin(\theta+\alpha) - v_m\cos\theta\right).$$

\par\smallskip\noindent\textbf{Answer:} \textbackslash{}textbf{Answer:} $$v_2 = \sqrt{v_1^2+v_m^2-2v_mv_1 s'}$$
\end{solcard}

\begin{solcard}{B3}{By writing the ZSE in the Moon's CO for the position of the ship closest to the Moon and the position immediately after entering the attraction zone, obtain and solve the quadratic equation and obtain the exact expression for $v_1/v_m$. Express your answer in terms of $s$, $s'$, $c$, $b/r_c$ and $v_\text{II}^m/v_m$. Note. There is no need to select a root sign at this point. Do not forget that within the framework of the model used, the interaction with the Earth in this area does not need to be taken into account.}{1.20}

Note. There is no need to select a root sign at this point. Do not forget that within the framework of the model used, the interaction with the Earth in this area does not need to be taken into account.


Let the speed of the ship at the point closest to the Moon be $v$. Then from the law of conservation of momentum we obtain: $$vb = L_0^m.$$$$v = \frac{r_c}{b}\left(v_1s - v_mc\right)$$Energy conservation law: $$\frac{v^2}{2} - \frac{GM_m}{b} = \frac{v_2^2}{2} - \frac{GM_m}{r_c}.$$$$\frac{r_c^2}{b^2}\left(v_1s - v_mc\right)^2 - \left(v_\text{II}^m\right)^2 = v_1^2+v_m^2-2v_mv_1 s' - \left(v_\text{II}^m\right)^2\frac{b}{r_c}$$For удобства denote $Z = v_1/v_m$, $x = b/r_c$, $y = v_\text{II}^m/v_m$. $$\frac{1}{x^2}\left(Zs - c\right)^2 - y^2 = Z^2+1-2Z s' - y^2x$$$$Z^2 \left(s^2-x^2\right) - 2Z\left(sc-x^2s'\right)+c^2-x^2-y^2x^2(1-x) = 0$$$$Z = \frac{sc-x^2s'\pm\sqrt{(sc-x^2s')^2-(s^2-x^2)(c^2-x^2-y^2x^2(1-x))}}{s^2-x^2}$$

\par\smallskip\noindent\textbf{Answer:} \textbackslash{}textbf{Answer:} $$\dfrac{v_1}{v_m} = \frac{sc-\left(\dfrac{b}{r_c}\right)^2s'\pm\sqrt{\left(sc-\left(\dfrac{b}{r_c}\right)^2s'\right)^2-\left(s^2-\left(\dfrac{b}{r_c}\right)^2\right)\left(c^2-\left(\dfrac{b}{r_c}\right)^2-\left(\dfrac{v_{II}^m}{v_m}\right)^2\left(\dfrac{b}{r_c}\right)^2\left(1-\left(\dfrac{b}{r_c}\right)\right)\right)}}{s^2-\left(\dfrac{b}{r_c}\right)^2}$$
\end{solcard}

\begin{solcard}{B4}{Show that the expression for $v_1$ obtained in paragraph B4, when expanded to first order in powers of $b/r_c$, is reduced to the form: $$v_1 \approx \dfrac{v_m c}{s} + B_0\cdot \dfrac{b}{r_c}$$ Express $B_0$ via/through $v_m$, $v_\text{II}^m$, $s$, $c$ and $s'$. Выбирая sign корня, покажите, что $B_0 > 0$.}{1.00}

$$v_1 \approx \dfrac{v_m c}{s} + B_0\cdot \dfrac{b}{r_c}$$


Express $B_0$ in terms of $v_m$, $v_\text{II}^m$, $s$, $c$ and $s'$. When choosing a root sign, show that $B_0 > 0$.


We will expand to first order in $x$. $$Z = \frac{sc-x^2s'\pm\sqrt{(sc-x^2s')^2-(s^2-x^2)(c^2-x^2-y^2x^2(1-x))}}{s^2-x^2}$$Consider expression под корнем с точностью до второго порядка (он может дать вклад порядка первого порядка, if под корнем нет слагаемых нулевого and первого порядка, что выполняется в данном сrayае): $$(sc-x^2s')^2-(s^2-x^2)(c^2-x^2-y^2x^2(1-x)) \approx s^2c^2 - 2scs'x^2 - \left(s^2c^2-x^2\left(c^2+s^2+y^2s^2\right)\right)=\\=\left(c^2+s^2+y^2s^2-2scs'\right)x^2.$$Expression for $Z$ с точностью до первого порядка: $$Z = \frac{sc\pm x\sqrt{c^2+s^2+y^2s^2-2scs'}}{s^2}.$$$$v_1 = \frac{v_mc}{s}\pm\frac{v_m\sqrt{c^2+s^2+\left(\frac{v_\text{II}^m}{v_m}\right)^2s^2-2scs'}}{s^2}\cdot\frac{b}{r_c}$$It is easy to notice that without the second term we get the speed at which the ship flies exactly to the Moon. The "+" sign corresponds to the desired flight around the Moon, and the "-" sign corresponds to the flight from the other side.


Thus, the final expression for $v_1$ is: $$v_1 = \frac{v_mc}{s}+\frac{v_m\sqrt{c^2+s^2+\left(\frac{v_\text{II}^m}{v_m}\right)^2s^2-2scs'}}{s^2}\cdot\frac{b}{r_c}$$

\par\smallskip\noindent\textbf{Answer:} \textbackslash{}textbf{Answer:} $$B_0 = \frac{v_m\sqrt{c^2+s^2+\left(\frac{v_\text{II}^m}{v_m}\right)^2s^2-2scs'}}{s^2}$$
\end{solcard}

\begin{solcard}{B5}{Show that the expression for $v_1$, up to first order, is: $$v_1 \approx \dfrac{v_m}{\text{tg}~\theta} + A \alpha + B \frac{b}{r_c}$$ Express $A$ and $B$ via/through $v_m$, $v_\text{II}^m$, $\theta$.}{0.60}

$$v_1 \approx \dfrac{v_m}{\text{tg}~\theta} + A \alpha + B \frac{b}{r_c}$$


Express $A$ and $B$ in terms of $v_m$, $v_\text{II}^m$, $\theta$.


Note that in the answer of the previous paragraph it is possible to expand $B_0$ to zero order, since when taking into account the corrections we will obtain terms of a high order of smallness (at least of the second type $\alpha b/r_c$). $$B_0 \approx \frac{v_m\sqrt{\cos^2+\sin^2\theta+\left(\frac{v_\text{II}^m}{v_m}\right)^2\sin^2\theta}}{\sin^2\theta}=\frac{v_m\sqrt{1+\left(\frac{v_\text{II}^m}{v_m}\right)^2\sin^2\theta}}{\sin^2\theta}$$Expanding $v_mc/s$. $$\frac{v_mc}{s} = \frac{v_m \cos\theta}{\sin(\alpha+\theta)}\approx\frac{v_m \cos\theta}{\sin\theta+\alpha\cos\theta}\approx\frac{v_m}{\text{tg}\theta}\left(1-\frac{\alpha}{\text{tg}\theta}\right)$$$$v_1 = \frac{v_m}{\text{tg}\theta}-\frac{v_m}{\text{tg}^2\theta}\alpha+\frac{v_m\sqrt{1+\left(\frac{v_\text{II}^m}{v_m}\right)^2\sin^2\theta}}{\sin^2\theta}\cdot\frac{b}{r_c}$$

\par\smallskip\noindent\textbf{Answer:} \textbackslash{}textbf{Answer:} $$A = -\frac{v_m}{\text{tg}^2\theta}$$$$B = \frac{v_m\sqrt{1+\left(\frac{v_\text{II}^m}{v_m}\right)^2\sin^2\theta}}{\sin^2\theta}$$
\end{solcard}

\begin{solcard}{B6}{Let's return to the Earth's reference system to find an expression for the angle $\alpha$ between the speed $v_1$ and the straight line Earth - Moon. Find an expression for $\alpha$ between the speed $v_1$ and the straight line Earth - Moon up to first order in powers $r_c/r_m$ and $a/r_m$. Express your answer using $a$, $r_c$, $r_m$, $v_1$, $v_0$ and $\theta$. Note. Do not forget that within the framework of the model used, interaction with the Moon in this area does not need to be taken into account.}{0.60}

Find an expression for $\alpha$ between the speed $v_1$ and the straight line Earth - Moon up to first order in powers $r_c/r_m$ and $a/r_m$. Express your answer using $a$, $r_c$, $r_m$, $v_1$, $v_0$ and $\theta$.


Note. Do not forget that within the framework of the model used, interaction with the Moon in this area does not need to be taken into account.


Let's use the law of conservation of angular momentum relative to the Earth. Consider the projections $\vec v_1$ and the radius vector from the Earth $\vec r$: $$\vec v_1 = \begin{pmatrix}v_1\cos\alpha\\v_1\sin\alpha\end{pmatrix},$$$$\vec r = \begin{pmatrix}r_m - r_c\cos\theta\\r_c\sin\theta\end{pmatrix}.$$Momentum conservation law: $$v_0 a = v_1\sin\alpha(r_m-r_c\cos\theta) - v_1 r_c\cos\alpha \sin\theta.$$$$\frac{v_0 a}{v_1\sqrt{r _m^2-2r_mr_c\cos\theta +r_c^2}}=\sin\left(\alpha - \text{arctg}\left(\frac{r_c\sin\theta}{r_m-r_c\cos\theta}\right)\right)$$$$\alpha = \text{arctg}\left(\frac{r_c\sin\theta}{r_m-r_c\cos\theta}\right) + \arcsin\left(\frac{v_0 a}{v_1\sqrt{r _m^2-2r_mr_c\cos\theta +r_c^2}}\right)\approx\frac{r_c\sin\theta}{r_m}+\frac{v_0a}{v_1r_m}$$

\par\smallskip\noindent\textbf{Answer:} \textbackslash{}textbf{Answer:} With
\end{solcard}

\begin{solcard}{B7}{Writing the ZSE in the FR of the Earth and neglecting the terms of order $a/r_m$ and $r_c/r_m$, express the speed $v_0$ in terms of $v_1$ and $v_\text{II}^e$. Note. Do not forget that within the framework of the model used, interaction with the Moon in this area does not need to be taken into account.}{0.40}

Note. Do not forget that within the framework of the model used, interaction with the Moon in this area does not need to be taken into account.


Law of conservation of energy: $$\frac{v_0^2}{2}-\frac{GM_e}{a} = \frac{v_1^2}{2}.$$Note that potentialьную энергию в правой части уравнения не учитываем, since она имеет порядок $a/r_m$. $$v_0^2-\left(v_\text{II}^e\right)^2 = v_1^2$$

\par\smallskip\noindent\textbf{Answer:} \textbackslash{}textbf{Answer:} $$v_0 = \sqrt{\left(v_\text{II}^e\right)^2+v_1^2}$$
\end{solcard}

\begin{solcard}{B8}{By writing the equation of a conic section, you get an equation from which you can find $\theta$, through $v_1$, $\alpha$, $v_m$, $G$, $M_m$, $b$ and $r_c$.}{1.00}

Let us write down the equation of the canonical section $r=\dfrac{p}{1+e\cos{\varphi}}$. Let us write this equation for points whose distances are $r_c$ and $b$ $$b=\dfrac{p}{1+e},$$$$r_c=\dfrac{p}{1-e\cos{\varphi}}.$$Expressing from уравнений эксцентриситет and for/atравнивая we get $$\dfrac{p}{b}-1=\left(1-\dfrac{p}{r_c}\right)\cdot \dfrac{1}{\cos{\theta}}.$$Substituting $p$ по формуле $$p=\dfrac{\left(L_0^m m\right)^2}{GM_m m^2}=\dfrac{\left(v_1\sin(\alpha+\theta)-v_m\cos{\theta}\right)^2 r_c^2}{GM_m}$$

\par\smallskip\noindent\textbf{Answer:} \textbackslash{}textbf{Answer:} $$\dfrac{\left(v_1\sin(\alpha+\theta)-v_m\cos{\theta}\right)^2 r_c^2}{GM_mb}-1=\left(1-\dfrac{\left(v_1\sin(\alpha+\theta)-v_m\cos{\theta}\right)^2 r_c}{GM_m}\right)\cdot\dfrac{1}{\cos{\theta}}$$
\end{solcard}

\begin{solcard}{B9}{Find $\beta_2$ - the ratio of the initial mass of the rocket to the final mass when accelerating from speed $v_a$ - orbital speed when moving in a circular orbit of radius $a$ to speed $v_0$. Express your answer using $u$, $v_\text{II}^e$ and $v_a$. Calculate the numerical value $\beta_2$, take the quantities necessary for calculations from part A.}{0.40}

Calculate the numerical value of $\beta_2$, take the quantities necessary for calculations from part A.


Let's write the differential equation from point A1 without taking into account the force of attraction and resistance: $$mdv+udm=0,$$ $$\int \limits_{v_a}^{v_{II}^e} \dfrac{dv}{u}=-\int \limits_{m_0}^{m} \dfrac{dm}{m},$$ $$\dfrac{v_{II}^e-v_a}{u}=-\ln{\dfrac{m}{m_0}},$$ $$\beta_2=\exp{\left(\dfrac{v_{II}^e-v_a}{u}\right)}$$


$$mdv+udm=0,$$


$$\int \limits_{v_a}^{v_{II}^e} \dfrac{dv}{u}=-\int \limits_{m_0}^{m} \dfrac{dm}{m},$$


$$\dfrac{v_{II}^e-v_a}{u}=-\ln{\dfrac{m}{m_0}},$$


$$\beta_2=\exp{\left(\dfrac{v_{II}^e-v_a}{u}\right)}$$

\par\smallskip\noindent\textbf{Answer:} \textbackslash{}textbf{Answer:} $$\beta_2=\exp{\left(\dfrac{\sqrt{2}-1}{u}\cdot\sqrt{\dfrac{GM_e}{a}}\right)}=2.05$$
\end{solcard}

\begin{solcard}{C1}{Counting $v_2 = 1.5~\mathrm{km}/с$ (take the remaining necessary quantities from the previous parts of the problem), calculate numerically how many times the mass of the rocket will decrease during the landing process $\beta_3$.}{0.50}

Let's write down the ESE for moving the rocket from the beginning of the moon's gravitational zone to the point at which we enter an orbit close to the moon $$\dfrac{mv_2^2}{2}-\dfrac{GM_mm}{r_c}=\dfrac{mv^2}{2}-\dfrac{GM_mm}{b},$$ $$v=\sqrt{v_2^2-2GM_m\left(\dfrac{1}{r_c}-\dfrac{1}{b}\right)}.$$ Let us write формулу Циолковского, поrayенную в partе B9 s заменой signа скорости истечения gasов $$\dfrac{v}{u}=\ln{\dfrac{m_0}{m_3}}$$


$$\dfrac{v}{u}=\ln{\dfrac{m_0}{m_3}}$$

\par\smallskip\noindent\textbf{Answer:} \textbackslash{}textbf{Answer:} $$\beta_3=\exp{\left(\dfrac{\sqrt{v_2^2-2GM_m\left(\dfrac{1}{r_c}-\dfrac{1}{b}\right)}}{u}\right)}=1.83$$
\end{solcard}

\begin{solcard}{C2}{Calculate numerically what initial mass $m$ the rocket must have to deliver cargo weighing $40~т$ to the Moon.}{0.50}
\par\smallskip\noindent\textbf{Answer:} \textbackslash{}textbf{Answer:} $$m=\beta_1\beta_2\beta_3 m_{\text{gруз}}\approx2300~\text{т}$$
\end{solcard}

\end{document}